% !TeX root = main.tex

% ============================================================
% thesis-sync entry: 2026-07-19_inchworm-fabrication
% Type: design decision (fabrication process) — reinstates the paper's
%       commented-out "Materials and Manufacturing" subsection.
% Part of the INCHWORM ROBOT block (prior work, precedes the lizard sections).
% Insertable section block (no preamble). Requires: graphicx.
% ============================================================

\section{ACTUATOR FABRICATION}
\label{sec:inchworm-fabrication}

\subsection{Body Casting and Fiber Reinforcement}

The first step in the manufacturing process was fabricating the mold, shown
in Fig.~\ref{fig:inchworm-mold}, to cast the silicone body (EcoFlex 00-30,
Smooth-On). Three brass prisms with a fan-shaped cross-section were used to
form the inner cavities of the actuator's chambers. Two CNC-machined acrylic
pieces acted as spacers, keeping the brass cores separated by 0.5~mm to
define the internal wall thickness. An acrylic casing was used to keep the
mold's cover and spacers perfectly aligned.

\begin{figure}[tbh]
 \centering
  \includegraphics[height=32mm]{figures/mold.eps}
  \vspace*{-2mm}
  \caption{The actuator's silicone casting mold.}
  \label{fig:inchworm-mold}
\end{figure}

The second step consisted of wrapping inextensible fibers in parallel
circumferences, with a 1.5~mm spacing, around the actuator's external walls
after coating them in fresh silicone. As the final step for the main body, a
resin 3D-printed top cap (with pressure ports) and bottom cap were glued to
the silicone using KE-44T adhesive (Shin-Etsu Chemical). The final
dimensions of the assembled actuator are 66~mm in length and 18~mm in
diameter.

\subsection{Suction Cups and Soft Hinge}

The suction cups (Fig.~\ref{fig:inchworm-cup}) were fabricated with a rigid
central piece of Polyetherimide (PEI) to prevent the cup ceiling from
collapsing under the negative pressure required for adhesion. A compliant
silicone lip was cast around this rigid core to ensure the suction cup
creates an effective initial seal. Finally, a secondary silicone lip was
cast on the distal ends of the actuator to create the soft hinge, and the
suction cups were bonded to these lips using the KE-44T adhesive.

\begin{figure}[tbh]
 \centering
  \includegraphics[height=30mm]{figures/suctioncup.eps}
  \vspace*{-2mm}
  \caption{The robot's suction cups. (Left) Diagram. (Right) Fabricated
  suction cup.}
  \label{fig:inchworm-cup}
\end{figure}
